\section{Full derivation of the $d=3$ exclusion}

This chapter gives the complete derivation of the $d=3$ exclusion:
the surviving family, its state formula, its residue class, and the
finite-base contradiction.

\subsection{The segment equation}

For $d=3$, the word molecule is
\begin{equation*}
  \wordA(w)=25+5\cdot2^{u_1}+2^{u_1+u_2},
  \qquad
  W=u_1+u_2+u_3.
\end{equation*}
The segment equation is
\begin{equation*}
  2^{W+e-1}g+2^W\delta5^{j-1}
    =5^3g+25+5\cdot2^{u_1}+2^{u_1+u_2}.
\end{equation*}

\subsection{The surviving family}

The branch analysis reduces the parameters to
\begin{equation*}
  t_j=2,\qquad \delta=1,\qquad e=2,\qquad u_1=u_2=1,\qquad a=0,
\end{equation*}
which forces
\begin{equation*}
  u_3=1.
\end{equation*}
Then
\begin{equation*}
  \wordA(w)=25+10+4=39,
  \qquad
  W=3.
\end{equation*}
The segment equation becomes
\begin{equation*}
  2^{3+1}g+2^3\cdot5^{j-1}=5^3g+39,
\end{equation*}
i.e.
\begin{equation*}
  16g+8\cdot5^{j-1}=125g+39.
\end{equation*}
Rearranging,
\begin{equation*}
  109g=8\cdot5^{j-1}-39,
\end{equation*}
so
\begin{equation*}
  g=\frac{8\cdot5^{j-1}-39}{109}.
\end{equation*}

\subsection{The integrality condition}

The state $g$ is an integer, so
\begin{equation*}
  8\cdot5^{j-1}\equiv39\pmod{109}.
\end{equation*}
The inverse of $8$ modulo $109$ is $41$, because
\begin{equation*}
  8\cdot41=328=3\cdot109+1.
\end{equation*}
Therefore
\begin{equation*}
  5^{j-1}\equiv39\cdot41=1599\equiv73\pmod{109},
\end{equation*}
since $1599-14\cdot109=73$.

\subsection{The residue class}

The full family calculation combines this congruence with the
remaining branch conditions and produces
\begin{equation*}
  j-1\equiv923\pmod{1728}.
\end{equation*}
The modulus $1728$ is specific to the $d=3$ family calculation; the
paper records the final class and the consistency checks below.

\subsection{Consistency checks}

Repeated squaring modulo $109$ gives:
\begin{center}
\small
\begin{tabular}{@{}rrrrrrr@{}}
\toprule
$k$ & $5^k\bmod109$ & $k$ & $5^k\bmod109$ & $k$ & $5^k\bmod109$ \\
\midrule
1 & 5 & 16 & 89 & 256 & 26 \\
2 & 25 & 32 & 73 & 512 & 22 \\
4 & 80 & 64 & 97 & 923 & 73 \\
8 & 78 & 128 & 35 & & \\
\bottomrule
\end{tabular}
\end{center}
The entry $5^{923}\equiv73$ follows from
\begin{equation*}
  923=512+256+128+16+8+2+1
\end{equation*}
and the corresponding product of residues.  Finally,
\begin{equation*}
  8\cdot73=584\equiv39\pmod{109},
\end{equation*}
so the representative $923$ is consistent with the integrality of
$g$.

\subsection{The first big step}

For the surviving family, the segment weights are
\begin{equation*}
  u_1=u_2=u_3=1,
\end{equation*}
and the first step of weight at least $3$ occurs at depth $j+4$ with
weight $6$.  The finite base says that the first big step of the orbit
of $7$ occurs at depth $15$ and has weight exactly $3$.  Hence the
family contradicts the finite base.

\subsection{Lean statement}

The exclusion is compiled as
\texttt{d3\_\allowbreak family\_\allowbreak big\_\allowbreak
weight\_\allowbreak excluded}.  It uses the finite-prefix facts
\texttt{fullOrbit\_\allowbreak first\_\allowbreak t\_\allowbreak ge3\_
\allowbreak is\_\allowbreak exactly\_\allowbreak 3} and
\texttt{first\_\allowbreak big\_\allowbreak step\_\allowbreak unique}.

\subsection{Status}

\begin{center}
\small
\begin{tabular}{@{}p{0.56\textwidth}p{0.36\textwidth}@{}}
\toprule
Item & Status \\
\midrule
family reduction & math proven \\
state formula and integrality & math proven \\
residue class and consistency & math proven \\
finite-base contradiction & math proven \\
\texttt{d3\_\allowbreak family\_\allowbreak big\_\allowbreak
weight\_\allowbreak excluded} & Lean compiled \\
\bottomrule
\end{tabular}
\end{center}
